\clearpage
\section{Исследовательская часть}

\subsection{Класс входных данных}

Для параметризации использован один класс из трёх полных неориентированных графов по восемь вершин. Для каждого графа задан один и тот же бюджет $B=165$, время ребра лежит в диапазоне $[24;120]$ часов, стоимость --- в диапазоне $[10;50]$ условных денежных единиц. Оба края каждого диапазона присутствуют в каждом графе.

Метки синтетические, а имена вершин соответствуют портам. Первый граф описывает переходы между портами Тихоокеанского региона, второй --- Северной Атлантики, третий --- Индийского океана. Начальные значения генератора меток --- 601, 602 и 603 соответственно. Полные входные данные приведены в приложении~\ref{app:graphs}. В таблице~\ref{tab:reference} представлены эталоны полного перебора.

\begin{table}[H]\centering
\caption{Эталонные результаты на графах класса}
\label{tab:reference}
\begin{tabular}{@{}crrr@{}}
\toprule
Граф & Лимит $B$ & Время $L^*$, ч & Стоимость $C^*$ \\
\midrule
\tablerows{results/reference.tex}
\bottomrule
\end{tabular}
\end{table}

Ограничение активно для всех трёх графов. Если снять лимит, минимальные времена составят 281, 221 и 229 часов, а стоимости соответствующих маршрутов --- 210, 198 и 204. Рисунок~\ref{fig:ports} показывает один эталонный маршрут в графе $G_1$.

\diagram{ports}{Эталонный маршрут по портам графа $G_1$}{fig:ports}

\clearpage
\subsection{Параметризация}

Проверялись значения $\alpha\in\{0,1,2\}$, $\beta\in\{0,1,3\}$ и $\rho\in\{0{,}2;\,0{,}5;\,0{,}8\}$. Число итераций $D=32$, число обычных муравьёв $K=12$, число элитных муравьёв $e=4$. Константы $Q=100$ и $\varepsilon=0{,}05$ одинаковы во всех запусках.

Для каждой комбинации параметров и каждого графа выполнено $M=20$ запусков. В пределах графа одинаковые начальные состояния генератора использовались для всех комбинаций. Серии начинались со значений 10000, 11000 и 12000. Всего проверено \Configurations{} комбинаций и выполнено \TotalAntRuns{} запусков.

Для каждого запуска с найденным допустимым путём относительное отклонение рассчитывалось в процентах
\begin{equation}
\delta(P)=100\,\frac{L(P)-L^*}{L^*},
\end{equation}
где $L^*$ --- эталонное время для графа. При исчерпании попыток без маршрута принималось $\delta=+\infty$. По каждому графу рассчитаны максимум $\delta_{\max}$ и медиана $\widetilde\delta$ за 20 запусков. Основной критерий для класса --- максимум из трёх значений $\delta_{\max}$. При равенстве выбиралась меньшая медиана трёх графовых медиан. При сохранении равенства использовался лексикографический порядок $(\alpha,\beta,\rho)$.

\begin{table}[H]\centering
\caption{Относительное отклонение при выбранных параметрах}
\label{tab:winner}
\begin{tabular}{@{}crrr@{}}
\toprule
Граф & $\delta_{\max}$, \% & $\widetilde\delta$, \% & Попыток без результата \\
\midrule
\tablerows{results/winner.tex}
\bottomrule
\end{tabular}
\end{table}
В таблице~\ref{tab:winner} приведены результаты при $\alpha=\BestAlpha$, $\beta=\BestBeta$ и $\rho=\BestRho$. Для этой комбинации максимум по классу равен $\WorstWinner\,\%$, медиана трёх графовых медиан --- $\MedianWinner\,\%$. Все запуски завершились допустимым маршрутом для \AllWithoutMiss{} из \Configurations{} комбинаций. Наибольшее конечное отклонение среди проверенных конфигураций составило $\GridWorstFinite\,\%$. Полные результаты приведены в приложении~\ref{app:grid}.

\clearpage
\subsection{Сравнение времени работы и качества результата}

Для каждого $n$ от 5 до 11 взят отдельный полный граф с теми же диапазонами меток. Назначен бюджет $B=23(n-1)$. Графы сгенерированы с начальными значениями $1000+n$. Муравьиный метод запущен с выбранными на предыдущем этапе параметрами. Замеры проведены в Windows 11 (сборка 26200) на процессоре Intel Core i3-1115G4. Оба алгоритма реализованы в программе, собранной GCC 14.2.0 с параметрами \texttt{-std=c++17 -O2}.

Счётчик \texttt{std::chrono::steady\_clock} измерял длительность вызова функции поиска. Чтение файла, проверка исходных данных и форматирование вывода выполнялись вне области замера. Для каждого размера сначала выполнен пробный запуск в отдельном процессе, затем проведено по семь замеров каждого метода. В таблице~\ref{tab:bench} приведены медианы времени работы. Отклонения рассчитаны по семи запускам муравьиного алгоритма.

\begin{table}[H]\centering
\caption{Время работы и отклонение от точного решения}
\label{tab:bench}
{\small\setlength{\tabcolsep}{5pt}
\begin{tabular}{@{}rrrrrrrr@{}}
\toprule
$n$ & $B$ & $L^*$, ч & \multicolumn{2}{c}{Время, мкс} & \multicolumn{2}{c}{Отклонение, \%} & Пропуски\\
\cmidrule(lr){4-5}\cmidrule(lr){6-7}
 & & & перебор & муравьи & макс. & мед. & \\
\midrule
\tablerows{results/benchmark.tex}
\bottomrule
\end{tabular}}
\end{table}

В данном эксперименте муравьиный метод впервые оказался быстрее при $n=\FirstFasterSize$. На графе из 11 портов медианы составили $\LastExact$ мс у перебора и $\LastAnt$ мс у муравьиного метода, что даёт отношение около $\SpeedupLargest$. В семи сериях измерений зарегистрировано \BenchmarkMisses{} пропусков допустимого решения. Максимальное отклонение времени среди серий без пропусков составило $\BenchmarkMax\,\%$.

\clearpage
На рисунке~\ref{fig:timing} по логарифмической шкале виден быстрый рост времени перебора. При фиксированном числе попыток время муравьиного метода зависит от числа шагов в каждой построенной траектории.

\diagram{timing}{Зависимость медианного времени поиска от числа портов}{fig:timing}

\subsection{Рекомендации по настройке}

На графах исследованного класса начальная настройка $\alpha=\BestAlpha$, $\beta=\BestBeta$, $\rho=\BestRho$ при $D=32$ и $K=12$ дала наименьшее наблюдаемое максимальное отклонение. Такой же максимум получили \BestTies{} комбинации. При $\alpha=0$ феромон исключается из весов, при $\beta=0$ веса перестают напрямую учитывать время перехода. Результаты для всех режимов приведены в таблице~\ref{tab:grid}.

Для небольшого числа портов полный перебор позволяет найти точное решение или установить, что допустимых маршрутов нет. Муравьиный метод ограничивает число попыток. При переносе выбранных настроек на другие бюджеты и метки следует проверять качество маршрутов отдельно.

\subsection*{Вывод}

Параметризация охватила \TotalAntRuns{} запусков. По измерениям на отдельных графах с ростом числа портов преимущество муравьиного метода по времени увеличивалось. Точность эвристики зависела от графа и бюджета.
